% ---- Title page ----
\begin{titlepage}
\centering
{\scshape\large Bocconi University \par}
\vspace{4pt}
{\scshape Master of Science in Artificial Intelligence
\par}
\vspace{40pt}
{\large\itshape Master's Thesis \par}
\vspace{30pt}
{\Huge\bfseries Learnable, Task-Adaptive Structured Memory for LLM Agents \par}
\vspace{18pt}
{\large\itshape Learning \emph{what to store} and \emph{how to retrieve}:
a parameter vector $\thetavec$, optimized by black-box search,
across grid worlds and six real long-context benchmarks. \par}
\vspace{50pt}
{\large
\textbf{Candidate:} Stefan Uifalean \\[4pt]
\textbf{Student ID:} 3310360 \\[4pt]
\textbf{Supervisor:} Prof.\ Dirk Hovy \par}
\vfill
{\large Academic Year 2025--2026 \par}
\end{titlepage}

% ---- Declaration ----
\cleardoublepage
\thispagestyle{empty}
\section*{Declaration}
\noindent I declare that this thesis is my own original work and that all sources
and collaborations have been duly acknowledged. The code, data, and results
supporting it are available in the public repository cited in
Appendix~\ref{app:repro}.
\vspace{20pt}

\noindent Stefan Uifalean \\
Milan, June 2026

% ---- Abstract ----
\cleardoublepage
\thispagestyle{empty}
\section*{Abstract}
\noindent
Memory in LLM-based agents is almost always \emph{fixed}: a context window of a
given size, a RAG index with a fixed chunking and retrieval rule, or a vector
store that treats every event alike. Yet different tasks demand different memory
structures --- a multi-hop question needs entity relations, a navigation task
needs spatial transitions, an end-of-corpus question needs recency-independent
recall. This thesis asks whether an agent can \textbf{learn how to construct its
own memory} --- what to store, which concepts to track as nodes, and how to
score retrieval --- and whether that learned structure is \textbf{task-dependent}.

We parameterize memory construction by a vector $\thetavec$ (storage
probability, entity-creation threshold, temporal-edge probability, and learnable
retrieval weights \wgraph{}/\wembed{}/\wrec{}, up to ten dimensions in the most
complete system, \Vfour{}) and optimize $\thetavec$ --- not the policy, not the
value function --- by black-box search (Evolution Strategy, then CMA-ES). The
work proceeds in three stages of increasing realism. \textbf{Stage~1} (grid
proof-of-concept) shows the optimizer recovers \emph{different} $\thetavec$ for
different tasks and that adaptive $\thetavec$ improves over the fixed-memory
baseline (MultiHopKeyDoor success $2.5\% \rightarrow 27.5\%$). \textbf{Stage~2} scales to
a twelve-memory-system $\times$ four-environment benchmark with ablation,
zero-shot $\thetavec$-transfer, and sensitivity analysis; the ten-dimensional
\Vfour{} reaches the top performance cluster, and $\thetavec$ tuned for one
long-haystack task transfers to another but not across task families.
\textbf{Stage~3} moves to real LLMs (a \texttt{gpt-4o-mini} answerer)
doing corpus-cumulative question answering over \textbf{six} real long-context
benchmarks, with a corpus-mode CMA-ES re-tuning of $\thetavec$ and a
Claude cross-vendor judge.

The central Stage-3 finding, established after an adversarial self-audit, is
that \textbf{corpus-cumulative tuning shifts $\thetavec$ toward
recency-independent, embedding-driven retrieval, producing an end-of-corpus
accuracy lift that survives Holm-corrected held-out testing on two of the three
domain-coherent benchmarks} (FinanceBench, CUAD; QASPER is non-significant),
while two pre-registered domain-incoherent controls also lift and NarrativeQA is
undefined by construction. The tuning objective is a validated proxy for answer
quality (recall@8 vs.\ judge score, pooled Spearman $\rho=0.69$), and the
corpus-mode result is near-deterministic at temperature~0 (cross-seed
std.\ $0.012$). Held against fair, same-budget baselines the advantage is
\textbf{benchmark-dependent}: versus a corpus-tuned Okapi BM25 the learned memory
wins on FinanceBench, ties on QASPER, and loses on CUAD, and versus faithful
reimplementations of HippoRAG and MemGPT/Letta no single retrieval mechanism
dominates across corpora --- on CUAD the memory-system ranking even
\emph{reverses} between $50$ and the full $510$ contracts. These conclusions hold
at full corpus scale (all $510$ CUAD contracts, all $281$ QASPER papers), and the
per-task optimum is itself largely predictable --- a leave-one-benchmark-out model
recovers $0.80$ of the tuning lift from cheap, LLM-free corpus statistics. A
full-context
``dump-all'' baseline is \emph{statistically tied} with selective retrieval on
accuracy but roughly 18$\times$ more expensive --- so the case for learned,
task-adaptive memory at this scale is \textbf{efficiency and scalability}, not an
accuracy cliff.
